\begin{abstract}
    We study a crucial yet often overlooked issue inherent to Vision Transformers (ViTs): feature maps of these models exhibit grid-like artifacts (``Original features'' in \cref{fig:teaser}), which hurt the performance of ViTs in downstream dense prediction tasks such as semantic segmentation, depth prediction, and object discovery.
    We trace this issue down to the positional embeddings at the input stage.
    To mitigate this, we propose a two-stage denoising approach, termed \ourslong (\ours).
    In the first stage, we separate the clean features from those contaminated by positional artifacts by enforcing cross-view feature consistency with neural fields on a per-image basis. This per-image optimization process extracts artifact-free features from raw ViT outputs, providing clean feature estimates for offline applications. In the second stage, we train a lightweight transformer block to predict clean features from raw ViT outputs, leveraging the derived estimates of the clean features as supervision. Our method, \ours{}, does not require re-training the existing pre-trained ViTs, and is immediately applicable to any Vision Transformer architecture. We evaluate our method on a variety of representative ViTs (DINO, DeiT-III, EVA02, CLIP, DINOv2, DINOv2-reg) and demonstrate that \ours consistently improves existing state-of-the-art general-purpose models in semantic and geometric tasks across multiple datasets (\cref{fig:teaser}, right, \cref{tab:dense_results,tab:obj_det,tab:obj_discovery}).
    We hope our study will encourage a re-evaluation of ViT design, especially regarding the naive use of positional embeddings. Our code and checkpoints are publicly available in our \href{https://jiawei-yang.github.io/DenoisingViT/}{project page}.
\end{abstract}